% Nouveau EGA V (cote 25): archive cote 25, site page 68 (PDF page 69).
% typescript. Draft compared: none.
% AI-transcribed by Claude Opus 5.5 (claude-opus-5-5, adaptive thinking, effort high), Message Batches API, batch ega25b (msgbatch_01A7n5YuXubUPNK1zsepi8RG), 2026-11-13; checked against both supplied scan parts; see CHECK_DECISIONS.tsv.
\scanpage{25/68}
\mspage{63}

en particulier $\underline{X}^{(m)}$ est lisse sur X, avec des\hand{/}fibres géométriquement
intègres (et en fait des variétés rationelles)\var{\inserted[typed]{de dimension m(r-m),}}{ de dimension m(r-m),}{insertion}{ega25_25_068-v0001}. Bien entendu, la
vérification se réduit au c as X=P, et à ce titre appartient, tout
\var{\margin{|}}{\placedmargin{|}}{marginal note placed}{ega25_25_068-v0002}comme les considérations précédentes, au n$^{\text{o}}$ des généralités sur les
grassmani\var{\xout{z}}{}{typewriter strike-out}{ega25_25_068-v0003}ennes (qui , je le vois déjà, vont \var{\over{\unsure{f}}{f}}{f}{writing over}{ega25_25_068-v0004}inir par être ma\var{\over{\unsure{gu}}{\unsure{g}}}{\unsure{g}}{writing over}{ega25_25_068-v0005}ifiés
par toi en un paragraphe séparé). Nous sommes maintenant en posses-%
sion d 'un analogue parfait du diagramme du n$^{\text{o}}$1 . Encore un point
oublié: comme préschéma au dessus de Gr$^{m}$ , $\underline{H}^{(m)}$ est canoniquement
isomorphe à P$(\underline{E}_{\mathrm{Gr}}\ /\overset{\hand{\vee}}{\underline{G}}\ )$, c'est donc un excellent fibré projectif
(mais on ne peut rien en conclure bien sur en général pour $\underline{X}^{(m)}$ sur
Gr$^{m}$).

La proposition 2.1. se transpose sans changement. Dans la
propositi on 2.2. , il faut lire: il f et s que pour tout $x \in Z$, on
ait \var{\xout{dimxx}}{}{typewriter strike-out}{ega25_25_068-v0006} dim $\overline{x} \hand{\leqslant}$ m-\hand{1}. Pour la démonstration, on peut par exemple
se ramener à 2.6. en considérant une \var{\xout{section}}{}{typewriter strike-out}{ega25_25_068-v0007} variété linéaire de
codimension m générique comme intersection de m hyperplans génériques
indépendants, Di\unsure{eu|uu}donné demerdetur. Du point de vue rédaction, si
(comme il me s emble préférable) on fa\var{\xout{i}}{}{typewriter strike-out}{ega25_25_068-v0008}it dès le début m général, \unsure{il}
semble préférable de prouver en même temps 2.6. , où bien entendu
dim X -1 est remplacé par dim X -m (en sous-entendant qu'une dimension
$\hand{<}$ 0 dans la formule si\var{\xout{\unsure{x}}}{}{typewriter strike-out}{ega25_25_068-v0009}gnifie que l'ensemble considéré est vide).
Corollaire 2.3. se lit en r emplaçant ``fini'' par ``de dimension $\hand{\leqslant}$ \var{\xout{mxx}}{}{typewriter strike-out}{ega25_25_068-v0010}
m-1 '' . Corollaire 2.4. pareil ; De même pour 2.5. , remplacer
dim f(X$_{i}$ ) $\hand{>}$ 0 par dim f(X$_{i}$) $\hand{\geqslant}$ m , et même changement dans 2.7.
La proposition 2.8.reste valable telle quelle, dans 2.9. remplacer
fini par dim $\hand{\leqslant}$m-1. Kif kif pour 2.10. , 2.11. . L'énoncé
2.12. reste valable tel quel, avec une démonstration essentiellement

\var{\margin{68}}{\placedmargin{68}}{marginal note placed}{ega25_25_068-v0011}
